% --- Archon physics formalization source begin ---
% archon:physics
% archon:covers PhyXMiniProblems/problem_phyx_mini_0034.lean
% archon:source-report reports/phyx_mini/problem_phyx_mini_0034.source.json

\chapter{Physics problem phyx\_mini\_0034}
\label{ch:PhyXMiniProblems_problem_phyx_mini_0034}

\paragraph{Problem source.}
\#\# Physical scenario

An observer to the right of the mirror-lens combination shown in figure (not to scale) sees two real images that are the same size and in the same location. One image is upright, and the other is inverted. Both images are 1.50 times larger than the object. The lens has a focal length of \textbackslash{}( 10.0 \textbackslash{}text\{ cm\} \textbackslash{}). The lens and mirror are separated by \textbackslash{}( 40.0 \textbackslash{}text\{ cm\} \textbackslash{}).

\#\# Question

Determine the focal length of the mirror.

\#\# Answer choices

- A: A: \textbackslash{}( +11.1 \textbackslash{}text\{ cm\} \textbackslash{})

- B: B: \textbackslash{}( +10.5 \textbackslash{}text\{ cm\} \textbackslash{})

- C: C: \textbackslash{}( +8.7 \textbackslash{}text\{ cm\} \textbackslash{})

- D: D: \textbackslash{}( +11.7 \textbackslash{}text\{ cm\} \textbackslash{})

\#\# Figure caption (auxiliary; use the image as primary evidence)

The image depicts a diagram involving optical elements:

1. **Mirror**: On the left side, there is a concave mirror labeled "Mirror." It reflects the object to the lens.

2. **Object**: An arrow labeled "Object" is positioned upright in front of the mirror.

3. **Lens**: In the middle, there is a converging (convex) lens labeled "Lens." 

4. **Images**: To the right of the lens, there is a vertical red arrow labeled "Images," indicating the lens's role in forming an image from the object.

5. **Eye**: On the far right, there is a depiction of an eye looking towards the lens.

The diagram illustrates the process of an image being formed by an optical system consisting of a mirror and a lens, observed by an eye.

\#\# Dataset metadata

Geometrical Optics; Physical Model Grounding Reasoning, Multi-Formula Reasoning

\paragraph{Recorded answer/context.}
D: D: \textbackslash{}( +11.7 \textbackslash{}text\{ cm\} \textbackslash{})

\paragraph{Figure/image path.}
/root/proposal\_for\_physic/science-mango/phyx\_mini\_run/phyx\_data/test\_image/34.png

\paragraph{Formalization target.}
create a compiling Lean file with sorry bodies at `PhyXMiniProblems/problem\_phyx\_mini\_0034.lean`. The Lean declarations must preserve the physical quantities, dimensions or dimensional roles, figure labels, governing-law hypotheses, and final relation expressed by this problem.
Use Mathlib/Physlib names found through LeanExplore where available. If a physics API is missing, introduce faithful local abstractions rather than scalar placeholder aliases.

\begin{theorem}[Physics formalization target]
\label{thm:physics:phyx_mini_0034:target}
\lean{PhyXMiniProblems.ProblemPhyXMini0034.problem_phyx_mini_0034}
\uses{def:physics:phyx-mini-0034:phyxminiproblems-problemphyxmini0034-centimetersvalue, def:physics:phyx-mini-0034:phyxminiproblems-problemphyxmini0034-mirrorlenssetup, def:physics:phyx-mini-0034:phyxminiproblems-problemphyxmini0034-directlensobjectdistancecm, def:physics:phyx-mini-0034:phyxminiproblems-problemphyxmini0034-directlensimagedistancecm, def:physics:phyx-mini-0034:phyxminiproblems-problemphyxmini0034-mirrorobjectdistancecm, def:physics:phyx-mini-0034:phyxminiproblems-problemphyxmini0034-mirrorimagedistancecm, def:physics:phyx-mini-0034:phyxminiproblems-problemphyxmini0034-reflectedlensobjectdistancecm, def:physics:phyx-mini-0034:phyxminiproblems-problemphyxmini0034-reflectedlensimagedistancecm, def:physics:phyx-mini-0034:phyxminiproblems-problemphyxmini0034-matchesfigureandgivenlengths, def:physics:phyx-mini-0034:phyxminiproblems-problemphyxmini0034-hasobservedcoincidentimages, def:physics:phyx-mini-0034:phyxminiproblems-problemphyxmini0034-gaussianimaginglawcm, def:physics:phyx-mini-0034:phyxminiproblems-problemphyxmini0034-transversemagnificationlawcm, def:physics:phyx-mini-0034:phyxminiproblems-problemphyxmini0034-twostagemagnificationlaw, def:physics:phyx-mini-0034:phyxminiproblems-problemphyxmini0034-matchesanswertonearesttenth}
The assigned autoformalize agent should translate this physics problem into Lean declarations in the covered file, with theorem and lemma proof bodies written as `by sorry`.
\end{theorem}
\begin{proof}
This is an autoformalization task, not a proof task. Produce statements that can later be proved by the physics prover without weakening the physical model.
\end{proof}

% --- Archon declaration topology begin ---
\section{Declaration topology}
\begin{definition}[Optical Length]
  \label{def:physics:phyx-mini-0034:phyxminiproblems-problemphyxmini0034-opticallength}
  \lean{PhyXMiniProblems.ProblemPhyXMini0034.OpticalLength}
  A signed physical length whose value transforms coherently with unit choice
\end{definition}
\begin{proof}
  This is the declared model definition of Optical Length.
\end{proof}
\begin{definition}[length In]
  \label{def:physics:phyx-mini-0034:phyxminiproblems-problemphyxmini0034-lengthin}
  \lean{PhyXMiniProblems.ProblemPhyXMini0034.lengthIn}
  \uses{def:physics:phyx-mini-0034:phyxminiproblems-problemphyxmini0034-opticallength}
  The physical length with scalar readout magnitude in the selected unit
\end{definition}
\begin{proof}
  This is the declared model definition of length In.
\end{proof}
\begin{definition}[centimeters Value]
  \label{def:physics:phyx-mini-0034:phyxminiproblems-problemphyxmini0034-centimetersvalue}
  \lean{PhyXMiniProblems.ProblemPhyXMini0034.centimetersValue}
  \uses{def:physics:phyx-mini-0034:phyxminiproblems-problemphyxmini0034-opticallength}
  The signed scalar readout of a physical length in centimeters
\end{definition}
\begin{proof}
  This is the declared model definition of centimeters Value.
\end{proof}
\begin{definition}[axial Separation Cm]
  \label{def:physics:phyx-mini-0034:phyxminiproblems-problemphyxmini0034-axialseparationcm}
  \lean{PhyXMiniProblems.ProblemPhyXMini0034.axialSeparationCm}
  \uses{def:physics:phyx-mini-0034:phyxminiproblems-problemphyxmini0034-opticallength, def:physics:phyx-mini-0034:phyxminiproblems-problemphyxmini0034-centimetersvalue}
  Directed axial separation, in centimeters, from start to finish
\end{definition}
\begin{proof}
  This is the declared model definition of axial Separation Cm.
\end{proof}
\begin{definition}[Thin Lens Kind]
  \label{def:physics:phyx-mini-0034:phyxminiproblems-problemphyxmini0034-thinlenskind}
  \lean{PhyXMiniProblems.ProblemPhyXMini0034.ThinLensKind}
  The paraxial lens types, distinguished by the sign of focal length
\end{definition}
\begin{proof}
  This is the declared model definition of Thin Lens Kind.
\end{proof}
\begin{definition}[Spherical Mirror Kind]
  \label{def:physics:phyx-mini-0034:phyxminiproblems-problemphyxmini0034-sphericalmirrorkind}
  \lean{PhyXMiniProblems.ProblemPhyXMini0034.SphericalMirrorKind}
  The spherical-mirror types as viewed from the incident-light side
\end{definition}
\begin{proof}
  This is the declared model definition of Spherical Mirror Kind.
\end{proof}
\begin{definition}[Image Orientation]
  \label{def:physics:phyx-mini-0034:phyxminiproblems-problemphyxmini0034-imageorientation}
  \lean{PhyXMiniProblems.ProblemPhyXMini0034.ImageOrientation}
  Orientation of an image relative to the upright object
\end{definition}
\begin{proof}
  This is the declared model definition of Image Orientation.
\end{proof}
\begin{definition}[Image Reality]
  \label{def:physics:phyx-mini-0034:phyxminiproblems-problemphyxmini0034-imagereality}
  \lean{PhyXMiniProblems.ProblemPhyXMini0034.ImageReality}
  Whether geometrical-optics rays actually meet at an image or only appear to
\end{definition}
\begin{proof}
  This is the declared model definition of Image Reality.
\end{proof}
\begin{definition}[Mirror Lens Setup]
  \label{def:physics:phyx-mini-0034:phyxminiproblems-problemphyxmini0034-mirrorlenssetup}
  \lean{PhyXMiniProblems.ProblemPhyXMini0034.MirrorLensSetup}
  \uses{def:physics:phyx-mini-0034:phyxminiproblems-problemphyxmini0034-opticallength, def:physics:phyx-mini-0034:phyxminiproblems-problemphyxmini0034-thinlenskind, def:physics:phyx-mini-0034:phyxminiproblems-problemphyxmini0034-sphericalmirrorkind, def:physics:phyx-mini-0034:phyxminiproblems-problemphyxmini0034-imageorientation, def:physics:phyx-mini-0034:phyxminiproblems-problemphyxmini0034-imagereality}
  The physical elements, labeled image locations, and path magnifications in the mirror--lens diagram. mirrorIntermediateImageAxialPosition is the image that the mirror alone would form. It is the effective object location for the reflected path's subsequent traversal of the lens; it is not assumed to coincide with the original object
\end{definition}
\begin{proof}
  This is the declared model definition of Mirror Lens Setup.
\end{proof}
\begin{definition}[direct Lens Object Distance Cm]
  \label{def:physics:phyx-mini-0034:phyxminiproblems-problemphyxmini0034-directlensobjectdistancecm}
  \lean{PhyXMiniProblems.ProblemPhyXMini0034.directLensObjectDistanceCm}
  \uses{def:physics:phyx-mini-0034:phyxminiproblems-problemphyxmini0034-axialseparationcm, def:physics:phyx-mini-0034:phyxminiproblems-problemphyxmini0034-mirrorlenssetup}
  Object distance for the direct lens path, as a positive centimeter readout
\end{definition}
\begin{proof}
  This is the declared model definition of direct Lens Object Distance Cm.
\end{proof}
\begin{definition}[direct Lens Image Distance Cm]
  \label{def:physics:phyx-mini-0034:phyxminiproblems-problemphyxmini0034-directlensimagedistancecm}
  \lean{PhyXMiniProblems.ProblemPhyXMini0034.directLensImageDistanceCm}
  \uses{def:physics:phyx-mini-0034:phyxminiproblems-problemphyxmini0034-axialseparationcm, def:physics:phyx-mini-0034:phyxminiproblems-problemphyxmini0034-mirrorlenssetup}
  Image distance for the direct lens path, in centimeters
\end{definition}
\begin{proof}
  This is the declared model definition of direct Lens Image Distance Cm.
\end{proof}
\begin{definition}[mirror Object Distance Cm]
  \label{def:physics:phyx-mini-0034:phyxminiproblems-problemphyxmini0034-mirrorobjectdistancecm}
  \lean{PhyXMiniProblems.ProblemPhyXMini0034.mirrorObjectDistanceCm}
  \uses{def:physics:phyx-mini-0034:phyxminiproblems-problemphyxmini0034-axialseparationcm, def:physics:phyx-mini-0034:phyxminiproblems-problemphyxmini0034-mirrorlenssetup}
  Object distance at the mirror, measured from mirror to object in centimeters
\end{definition}
\begin{proof}
  This is the declared model definition of mirror Object Distance Cm.
\end{proof}
\begin{definition}[mirror Image Distance Cm]
  \label{def:physics:phyx-mini-0034:phyxminiproblems-problemphyxmini0034-mirrorimagedistancecm}
  \lean{PhyXMiniProblems.ProblemPhyXMini0034.mirrorImageDistanceCm}
  \uses{def:physics:phyx-mini-0034:phyxminiproblems-problemphyxmini0034-axialseparationcm, def:physics:phyx-mini-0034:phyxminiproblems-problemphyxmini0034-mirrorlenssetup}
  Image distance at the mirror, measured into the reflected-light side
\end{definition}
\begin{proof}
  This is the declared model definition of mirror Image Distance Cm.
\end{proof}
\begin{definition}[reflected Lens Object Distance Cm]
  \label{def:physics:phyx-mini-0034:phyxminiproblems-problemphyxmini0034-reflectedlensobjectdistancecm}
  \lean{PhyXMiniProblems.ProblemPhyXMini0034.reflectedLensObjectDistanceCm}
  \uses{def:physics:phyx-mini-0034:phyxminiproblems-problemphyxmini0034-axialseparationcm, def:physics:phyx-mini-0034:phyxminiproblems-problemphyxmini0034-mirrorlenssetup}
  Object distance for the lens encounter after reflection. The mirror's intermediate image is the effective object for light incident on the lens from the left
\end{definition}
\begin{proof}
  This is the declared model definition of reflected Lens Object Distance Cm.
\end{proof}
\begin{definition}[reflected Lens Image Distance Cm]
  \label{def:physics:phyx-mini-0034:phyxminiproblems-problemphyxmini0034-reflectedlensimagedistancecm}
  \lean{PhyXMiniProblems.ProblemPhyXMini0034.reflectedLensImageDistanceCm}
  \uses{def:physics:phyx-mini-0034:phyxminiproblems-problemphyxmini0034-axialseparationcm, def:physics:phyx-mini-0034:phyxminiproblems-problemphyxmini0034-mirrorlenssetup}
  Image distance for the lens encounter after reflection, in centimeters
\end{definition}
\begin{proof}
  This is the declared model definition of reflected Lens Image Distance Cm.
\end{proof}
\begin{definition}[Matches Figure And Given Lengths]
  \label{def:physics:phyx-mini-0034:phyxminiproblems-problemphyxmini0034-matchesfigureandgivenlengths}
  \lean{PhyXMiniProblems.ProblemPhyXMini0034.MatchesFigureAndGivenLengths}
  \uses{def:physics:phyx-mini-0034:phyxminiproblems-problemphyxmini0034-centimetersvalue, def:physics:phyx-mini-0034:phyxminiproblems-problemphyxmini0034-axialseparationcm, def:physics:phyx-mini-0034:phyxminiproblems-problemphyxmini0034-mirrorlenssetup}
  Figure and numerical readouts: a concave mirror lies left of the upright object, which lies left of a converging lens; both displayed images lie between the lens and the eye. The lens focal length is +10.0 cm and the mirror--lens separation is 40.0 cm
\end{definition}
\begin{proof}
  This is the declared model definition of Matches Figure And Given Lengths.
\end{proof}
\begin{definition}[Has Observed Coincident Images]
  \label{def:physics:phyx-mini-0034:phyxminiproblems-problemphyxmini0034-hasobservedcoincidentimages}
  \lean{PhyXMiniProblems.ProblemPhyXMini0034.HasObservedCoincidentImages}
  \uses{def:physics:phyx-mini-0034:phyxminiproblems-problemphyxmini0034-mirrorlenssetup}
  Observed image data: the two real images occupy the same location, the direct image is inverted, the reflected-path image is upright, and their signed magnifications are respectively -1.50 and +1.50
\end{definition}
\begin{proof}
  This is the declared model definition of Has Observed Coincident Images.
\end{proof}
\begin{definition}[Gaussian Imaging Law Cm]
  \label{def:physics:phyx-mini-0034:phyxminiproblems-problemphyxmini0034-gaussianimaginglawcm}
  \lean{PhyXMiniProblems.ProblemPhyXMini0034.GaussianImagingLawCm}
  \uses{def:physics:phyx-mini-0034:phyxminiproblems-problemphyxmini0034-opticallength, def:physics:phyx-mini-0034:phyxminiproblems-problemphyxmini0034-centimetersvalue}
  The paraxial Gaussian imaging law 1/f = 1/p + 1/q, in the division-free, dimensionally homogeneous form f * (p + q) = p * q. All three arguments are centimeter readouts of signed lengths
\end{definition}
\begin{proof}
  This is the declared model definition of Gaussian Imaging Law Cm.
\end{proof}
\begin{definition}[Transverse Magnification Law Cm]
  \label{def:physics:phyx-mini-0034:phyxminiproblems-problemphyxmini0034-transversemagnificationlawcm}
  \lean{PhyXMiniProblems.ProblemPhyXMini0034.TransverseMagnificationLawCm}
  The signed transverse-magnification law m = -q/p, without division
\end{definition}
\begin{proof}
  This is the declared model definition of Transverse Magnification Law Cm.
\end{proof}
\begin{definition}[Two Stage Magnification Law]
  \label{def:physics:phyx-mini-0034:phyxminiproblems-problemphyxmini0034-twostagemagnificationlaw}
  \lean{PhyXMiniProblems.ProblemPhyXMini0034.TwoStageMagnificationLaw}
  \uses{def:physics:phyx-mini-0034:phyxminiproblems-problemphyxmini0034-mirrorlenssetup}
  The net magnification of the mirror and subsequent lens is their product
\end{definition}
\begin{proof}
  This is the declared model definition of Two Stage Magnification Law.
\end{proof}
\begin{lemma}[direct Path Distances eq]
  \label{lem:physics:phyx-mini-0034:phyxminiproblems-problemphyxmini0034-directpathdistances-eq}
  \lean{PhyXMiniProblems.ProblemPhyXMini0034.directPathDistances_eq}
  \uses{def:physics:phyx-mini-0034:phyxminiproblems-problemphyxmini0034-mirrorlenssetup, def:physics:phyx-mini-0034:phyxminiproblems-problemphyxmini0034-directlensobjectdistancecm, def:physics:phyx-mini-0034:phyxminiproblems-problemphyxmini0034-directlensimagedistancecm, def:physics:phyx-mini-0034:phyxminiproblems-problemphyxmini0034-matchesfigureandgivenlengths, def:physics:phyx-mini-0034:phyxminiproblems-problemphyxmini0034-hasobservedcoincidentimages, def:physics:phyx-mini-0034:phyxminiproblems-problemphyxmini0034-gaussianimaginglawcm, def:physics:phyx-mini-0034:phyxminiproblems-problemphyxmini0034-transversemagnificationlawcm}
  The direct path fixes the lens object and image distances
\end{lemma}
\begin{proof}
  The conclusion follows from the governing relations and figure readouts named in the statement of direct Path Distances eq.
\end{proof}
\begin{lemma}[mirror Intermediate Image coincides with object]
  \label{lem:physics:phyx-mini-0034:phyxminiproblems-problemphyxmini0034-mirrorintermediateimage-coincides-with-object}
  \lean{PhyXMiniProblems.ProblemPhyXMini0034.mirrorIntermediateImage_coincides_with_object}
  \uses{def:physics:phyx-mini-0034:phyxminiproblems-problemphyxmini0034-mirrorlenssetup, def:physics:phyx-mini-0034:phyxminiproblems-problemphyxmini0034-directlensobjectdistancecm, def:physics:phyx-mini-0034:phyxminiproblems-problemphyxmini0034-directlensimagedistancecm, def:physics:phyx-mini-0034:phyxminiproblems-problemphyxmini0034-reflectedlensobjectdistancecm, def:physics:phyx-mini-0034:phyxminiproblems-problemphyxmini0034-reflectedlensimagedistancecm, def:physics:phyx-mini-0034:phyxminiproblems-problemphyxmini0034-matchesfigureandgivenlengths, def:physics:phyx-mini-0034:phyxminiproblems-problemphyxmini0034-hasobservedcoincidentimages, def:physics:phyx-mini-0034:phyxminiproblems-problemphyxmini0034-gaussianimaginglawcm, def:physics:phyx-mini-0034:phyxminiproblems-problemphyxmini0034-transversemagnificationlawcm}
  Because the same lens sends both paths to the same final image plane, the mirror's intermediate image must coincide axially with the original object
\end{lemma}
\begin{proof}
  The conclusion follows from the governing relations and figure readouts named in the statement of mirror Intermediate Image coincides with object.
\end{proof}
\begin{lemma}[mirror Focal Length eq]
  \label{lem:physics:phyx-mini-0034:phyxminiproblems-problemphyxmini0034-mirrorfocallength-eq}
  \lean{PhyXMiniProblems.ProblemPhyXMini0034.mirrorFocalLength_eq}
  \uses{def:physics:phyx-mini-0034:phyxminiproblems-problemphyxmini0034-centimetersvalue, def:physics:phyx-mini-0034:phyxminiproblems-problemphyxmini0034-mirrorlenssetup, def:physics:phyx-mini-0034:phyxminiproblems-problemphyxmini0034-directlensobjectdistancecm, def:physics:phyx-mini-0034:phyxminiproblems-problemphyxmini0034-directlensimagedistancecm, def:physics:phyx-mini-0034:phyxminiproblems-problemphyxmini0034-mirrorobjectdistancecm, def:physics:phyx-mini-0034:phyxminiproblems-problemphyxmini0034-mirrorimagedistancecm, def:physics:phyx-mini-0034:phyxminiproblems-problemphyxmini0034-reflectedlensobjectdistancecm, def:physics:phyx-mini-0034:phyxminiproblems-problemphyxmini0034-reflectedlensimagedistancecm, def:physics:phyx-mini-0034:phyxminiproblems-problemphyxmini0034-matchesfigureandgivenlengths, def:physics:phyx-mini-0034:phyxminiproblems-problemphyxmini0034-hasobservedcoincidentimages, def:physics:phyx-mini-0034:phyxminiproblems-problemphyxmini0034-gaussianimaginglawcm, def:physics:phyx-mini-0034:phyxminiproblems-problemphyxmini0034-transversemagnificationlawcm}
  The supplied geometry and imaging laws determine the mirror focal length exactly as 35/3 cm
\end{lemma}
\begin{proof}
  The conclusion follows from the governing relations and figure readouts named in the statement of mirror Focal Length eq.
\end{proof}
\begin{definition}[Answer Choice]
  \label{def:physics:phyx-mini-0034:phyxminiproblems-problemphyxmini0034-answerchoice}
  \lean{PhyXMiniProblems.ProblemPhyXMini0034.AnswerChoice}
  Labels of the four printed multiple-choice answers
\end{definition}
\begin{proof}
  This is the declared model definition of Answer Choice.
\end{proof}
\begin{definition}[answer Focal Length Cm]
  \label{def:physics:phyx-mini-0034:phyxminiproblems-problemphyxmini0034-answerfocallengthcm}
  \lean{PhyXMiniProblems.ProblemPhyXMini0034.answerFocalLengthCm}
  \uses{def:physics:phyx-mini-0034:phyxminiproblems-problemphyxmini0034-answerchoice}
  Mirror focal-length readout printed beside each answer, in centimeters
\end{definition}
\begin{proof}
  This is the declared model definition of answer Focal Length Cm.
\end{proof}
\begin{definition}[Matches Answer To Nearest Tenth]
  \label{def:physics:phyx-mini-0034:phyxminiproblems-problemphyxmini0034-matchesanswertonearesttenth}
  \lean{PhyXMiniProblems.ProblemPhyXMini0034.MatchesAnswerToNearestTenth}
  \uses{def:physics:phyx-mini-0034:phyxminiproblems-problemphyxmini0034-opticallength, def:physics:phyx-mini-0034:phyxminiproblems-problemphyxmini0034-centimetersvalue, def:physics:phyx-mini-0034:phyxminiproblems-problemphyxmini0034-answerchoice, def:physics:phyx-mini-0034:phyxminiproblems-problemphyxmini0034-answerfocallengthcm}
  Agreement with a displayed focal length to the nearest tenth centimeter
\end{definition}
\begin{proof}
  This is the declared model definition of Matches Answer To Nearest Tenth.
\end{proof}
% --- Archon declaration topology end ---
% --- Archon physics formalization source end ---
